% Pursuing Stacks (cotes 134-2 to 134-8): archive cote 134-2, site page 56 (PDF page 57).
% appendix (letters to Breen). Draft compared: ps1.tex:1372-1389.
% AI-transcribed by Claude Opus 5.5 (claude-opus-5-5, adaptive thinking, effort high), Message Batches API, batch ps1 (msgbatch_01Rv5hCaNTtUt4JHx1yqoHPL), 2026-11-12; not yet checked by eye.
\scanpage{134-2/56}
\margin{\unsure{41}}

\hand{Villecun 17/19 July 1975}

Dear Larry,

\margin{(App.9)}I am happy to finish by receiving an echo to my long letter and even a
beginning to a constructive approach to a theor\over{em}{y} of the type I envisaged.
The construction which you propose for the notion of a non-strict n-category,
and of the nerve functor, has certainly the merit of existing, and of being
a first precise approach, but otherwise can be subject to some evident
criticism: it is very technical, unintuitive (yet at the level of 1-cat, etc,
and even of 2-cat, everything is so clear\hand{\unsure{-}} ``you just follow your nose\inserted{-tip}...'').
And finally the absence of a definition of a functor sending (semi-)simplicial
sets to n-groupoids. This functor corresponds to a geometric intuition so
clear that a theory which does not include it seems to me \deleted{inadequate.} \inserted{kind of a joke!} Perhaps
in trying to write down (like a sort of list of Christmas presents!) in a
complete and explicit enough way the notions which one would like to have at
ones disposal, and the relations (functor, equivalence, etc.) which should link
them, one would arrive finally at a kind of axiomatic description sufficiently
complete which should either give the key to an explicit \textit{ad hoc} construction,
or should permit \deleted{maybe} \inserted{at least} to enunciate and prove a theorem of existence and
uniqueness \inserted{(25)} for a theory of the required type. \margin{$\times$}

Otherwise, not having understood the idea of Segal in your last letter
(which I have generously sent to Illusie ...), I do not see how you define
the Picard n-categories - but this matters little. As far as ``strict'' Picard
n-categories are concerned, all I ask of them is that they finally form an
(n+1)-category (n+1)-equivalent to that of \inserted{chain} complexes \deleted{of chains} of length $n$ .
Agreed? I thank you for having rectified in my mind a big blunder, due to
my great ignorance of algebraic topology and homotopy - I was in fact of the
impression that H-spaces satisfying conditions of associativity and comm-%
utativity strict enough (say equivalent to an $\Omega^{i}X$ with $i$ arbitrarily large)
correspond to commutative topological groups (inspired by several analogies...\hand{)}
Thus I am entirely in agreement with your observations on p.5.

\mspage{56}
